% Chapter 21 -- Can Foundation Models Forecast Better Than Classical Methods?
% Requires: amsmath, amssymb, booktabs, array, tabularx, longtable, graphicx, tcolorbox (breakable, skins), hyperref or url.
% Defines the callout box only if the preamble has not.
\providecolor{fenavy}{HTML}{1B2A6B}
\providecolor{feshade}{HTML}{F3F5FA}
\makeatletter
\@ifundefined{callout}{%
  \newtcolorbox{callout}{breakable, enhanced,
    colback=white, colframe=gray!70, boxrule=0.5pt, arc=1mm,
    borderline west={2.5pt}{0pt}{gray!70},
    left=8pt, right=8pt, top=4pt, bottom=4pt}%
}{}
\makeatother

\providecolor{fenavy}{HTML}{1B2A6B}
\providecolor{feshade}{HTML}{F3F5FA}
\makeatletter
% Shaded box for optional technical material that general readers may skip.
\@ifundefined{technote}{%
  \newtcolorbox{technote}{breakable, enhanced,
    colback=feshade, colframe=fenavy!60, boxrule=0.5pt, arc=1mm,
    left=8pt, right=8pt, top=4pt, bottom=4pt}%
}{}
\makeatother

\chapter{Can Foundation Models Forecast Better Than Classical Methods?}
\label{ch:21}
\providecommand{\term}[1]{\emph{#1}}

Chatbots that write essays and answer questions have made artificial intelligence part of everyday life. The same technology is now being applied to numbers, with the promise that a single pretrained model can forecast almost any series, from sales and electricity demand to stock returns and inflation, without being built for the job. This chapter asks whether that promise holds for financial and economic data.

\term{Large language models}\index{large language models} learned to write by reading most of the internet. In 2024 the same recipe was applied to numbers, in the hope that one model trained on many series could forecast new ones without a separate model for each. Ansari et al.\ (2024) introduced Chronos\index{foundation models!Chronos}, which turns a time series into tokens, small pieces that the model reads the way a chatbot reads words, and trains a language-model architecture on them. Das et al.\ (2024) introduced TimesFM\index{foundation models!TimesFM}. These time-series \term{foundation models}\index{foundation models} are first trained on very large collections of series, from electricity demand and web traffic to weather and retail sales. They then forecast series they have never seen, with no further training. Their developers reported that, used this way, the models match or beat models fitted specifically to each dataset.

For finance and economics the promise is striking. One pretrained model, applied to returns, volatility\index{volatility} or inflation\index{inflation}, could replace the careful work of building and estimating a model for each series. But financial series are unlike most of the data these models learn from. Returns are close to unpredictable. Volatility has long memory and fat tails\index{fat tails}. Macroeconomic series are short and are often revised. Public financial data may also have been part of the \term{training set}\index{training, validation and test sets}, which can make a model look better on history than it is. Rahimikia, Ni and Wang (2025) tested foundation models on stock returns from many countries. Used off the shelf, the models did not beat simple benchmarks, although training them on financial data helped.

The stakes are high for forecasting practice in finance. If a pretrained model forecasts as well as a carefully built one, much of the craft of time-series econometrics becomes unnecessary. If not, the benchmarks of Chapters~\ref{ch:06} and~\ref{ch:07} remain the standard. 

The question of this chapter is: \textbf{can general-purpose time-series foundation models, used zero-shot, forecast financial and economic series more accurately than classical models?}

The lab in Section~\ref{sec:21-lab} repeats the zero-shot comparison of these papers on six free series from the Kenneth R. French Data Library and FRED, using the open model weights of Chronos and TimesFM. It checks whether foundation models, used out of the box, forecast better than classical methods.


\section{Foundations}
\label{sec:21-foundations}

\subsection{Classical benchmarks}

A forecasting method deserves attention only if it beats simple alternatives. Section~\ref{sec:09-benchmarks} described the standard ones. The historical mean is hard to beat for stock returns (Chapter~\ref{ch:06}). The random walk\index{random walk} is hard to beat for exchange rates\index{exchange rate} and interest-rate levels. This chapter adds richer classical models:
\begin{itemize}
  \item An \term{autoregression}, the AR(1)\index{autoregressive model} model extended to $p$ lags, regresses the series on its own past $p$ values:
  \begin{equation}
    y_{t+1} = c + \phi_1 y_t + \cdots + \phi_p y_{t-p+1} + \varepsilon_{t+1}
    \label{eq:21-ar}
  \end{equation}
  The number of lags can be chosen by an information criterion such as BIC, a score that rewards a good fit but penalizes extra parameters. \term{ARIMA}\index{ARIMA|textbf} models add differencing and moving-average terms.
  \item For volatility, the GARCH and HAR models of Chapter~\ref{ch:07} (Sections~\ref{sec:10-garch} and~\ref{sec:10-rv}; Corsi 2009) capture volatility's long memory with few parameters.
\end{itemize}

\subsection{Foundation models}
\label{sec:21-fm}

A time-series foundation model is a large neural network\index{neural network}, usually a \term{transformer}\index{transformer|textbf}\label{kt:transformer}, the network design behind modern chatbots, pretrained on a broad collection of series and then used \term{zero-shot}\index{foundation models!zero-shot forecasting|textbf}, without any training on the target series. In plain words, these models chop a series into pieces and learn to predict the next piece, much as a chatbot predicts the next word. The shaded box describes the two best-known designs.

\begin{technote}
\textbf{Technical note: two model designs.} Two design families matter here:

\begin{itemize}
  \item \term{Chronos}\index{foundation models!Chronos|textbf} (Ansari et al.\ 2024) scales each series and quantizes its values into a fixed vocabulary of tokens. It then trains a T5 language-model architecture to predict the next token, which yields a full probability distribution for the forecast. Models range from 20 million to 710 million parameters. Chronos-Bolt is a faster successor.
  \item \term{TimesFM}\index{foundation models!TimesFM|textbf} (Das et al.\ 2024) is a decoder-only transformer of about 200 million parameters. It reads the series in patches of consecutive points and outputs longer patches. It was pretrained on about 100 billion time points, mostly Google Trends search interest and Wikipedia page views, plus synthetic series.
\end{itemize}
\end{technote}

The pitch is breadth: a model that has seen trends, seasonality, mean reversion and regime shifts across many domains should recognize them in a new series.

\subsection{Contamination}
\label{sec:21-contam}

A forecast evaluation is valid only if the model has not seen the test data. Classical models are fitted on past data at each forecast date, so this is easy to guarantee. A foundation model's training data are fixed and often only partly documented. They may contain the test series, or close copies such as a stock index that tracks the target. The model may then partly be remembering rather than forecasting. The clean test uses data published after the model's \term{training cutoff}\index{training cutoff|textbf}.

\subsection{Evaluation}

Forecasts are compared with the walk-forward\index{walk-forward evaluation} design of Section~\ref{sec:09-walkforward}, which forecasting papers call \term{rolling origins}\index{walk-forward evaluation!rolling origin|textbf}. At each month $t$, every model forecasts month $t+1$ using only data up to $t$ (Figure~\ref{fig:21-rolling}). The tools of Chapter~\ref{ch:06} then apply. The mean squared error of Section~\ref{sec:09-foundations} is reported as a ratio to the benchmark's, so a ratio below 1 means the model is better. The Diebold--Mariano test of Section~\ref{sec:09-accuracy} asks whether the difference is larger than chance. When a model forecasts a whole distribution, coverage checks its \term{calibration}\index{calibration}, and the \term{pinball loss}\index{loss function!pinball|textbf} scores each forecast quantile\index{quantile}, that is, each point of a forecast range such as the level the outcome should stay below 90\% of the time. For volatility, the QLIKE loss of Section~\ref{sec:10-eval} is used, because it is robust to noise in the realized variance\index{variance} that stands in for true volatility (Patton 2011).

\begin{figure}[htbp]
\centering
\scalebox{0.9}{%
\begin{tikzpicture}[x=0.95cm, y=0.62cm]
  % Four forecast origins; data grow by one month at each origin
  \foreach \wfI/\wfE in {1/4.6,2/5.6,3/6.6,4/7.6} {
    \fill[fenavy!40] (0,-\wfI+0.3) rectangle (\wfE,-\wfI-0.3);
    \fill[fenavy]    (\wfE+0.08,-\wfI+0.3) rectangle (\wfE+0.88,-\wfI-0.3);
    \node[anchor=west] at (9.0,-\wfI)
      {origin \wfI: data to $t_\wfI$, forecast $t_\wfI + 1$};
  }
  \draw[-stealth] (0,-5) -- (8.3,-5) node[anchor=west] {time};
  % Legend
  \fill[fenavy!40] (0,-6.3) rectangle (0.6,-5.8);
  \node[anchor=west] at (0.7,-6.05) {data the models may use};
  \fill[fenavy]    (5.2,-6.3) rectangle (5.8,-5.8);
  \node[anchor=west] at (5.9,-6.05) {month being forecast};
\end{tikzpicture}%
}
\caption{A fair test lets every model see only the past. Rolling-origin evaluation. At each origin every model sees the same data up to that month and forecasts the next one; the origin then moves forward a month. No model ever sees the month it is forecasting.}
\label{fig:21-rolling}
\end{figure}

\textbf{Worked example.} Suppose an AR model's one-month inflation forecasts have an MSE of 0.040 (percentage points squared) over 300 months and a foundation model's 0.038.

\begin{itemize}
  \item The MSE ratio is 0.95, a 5\% gain.
  \item If the monthly loss differences have a standard deviation\index{standard deviation} of 0.020 and little autocorrelation\index{autocorrelation}, the Diebold--Mariano statistic is $0.002 / (0.020/\sqrt{300}) \approx 1.7$, not significant at 5\%.
  \item A gain of that size needs a long evaluation sample to be distinguished from luck, and a post-cutoff window of two years has only 24 observations (Figure~\ref{fig:21-dm}).
\end{itemize}

\begin{figure}[htbp]
\centering
\scalebox{0.9}{%
\begin{tikzpicture}[
  % Expected DM statistic grows with the square root of the sample: 0.1 x sqrt(T)
  declare function={dm(\t) = 0.1*sqrt(\t);}
]
\begin{axis}[
  width=0.8\linewidth, height=0.42\linewidth,
  axis lines=left, axis line style={-stealth},
  xmin=0, xmax=500, ymin=0, ymax=2.5,
  xtick={0,100,...,500}, ytick={0,1,2},
  xlabel={Months in the evaluation sample},
  ylabel={Expected DM statistic},
  samples=200, domain=1:500, clip=false,
]
  \addplot[fegray, dashed, domain=0:500, samples=2] {1.96};
  \addplot[fenavy, very thick] {dm(x)};
  \foreach \wfT in {24,300,384} {
    \edef\wfTmp{\noexpand\fill (axis cs:\wfT,{dm(\wfT)}) circle (2.4pt);}\wfTmp
  }
  \node[fegray, anchor=north] at (axis cs:120,1.92) {5\% critical value, 1.96};
  \node[anchor=south east] at (axis cs:375,2.02) {384 months};
  \node[anchor=north west] at (axis cs:305,{dm(300)-0.03}) {300 months: 1.73};
  \node[anchor=north west] at (axis cs:30,{dm(24)+0.05}) {24 months: 0.49};
\end{axis}
\end{tikzpicture}%
}
\caption{A small forecasting gain takes decades of monthly data to prove. The expected Diebold--Mariano statistic for the worked example's 5\% gain, as the evaluation sample grows. It reaches 1.96 only after about 384 months; a two-year post-cutoff window gives about 0.5.}
\label{fig:21-dm}
\end{figure}

\section{How it works}
\label{sec:21-how}

A zero-shot forecast takes three steps. First, the recent history of a series, up to a few hundred observations, is rescaled and passed to the pretrained network. Second, the network outputs a forecast for the next step or steps, often as a set of quantiles. Third, that output is used as it is. Nothing is estimated on the target series, and all the knowledge comes from \term{pretraining}.

That design has clear strengths:

\begin{itemize}
  \item \textbf{No specification work.} There are no lags to choose and no model to estimate.
  \item \textbf{Short histories.} Patterns learned elsewhere can fill in for a lack of data on the target series.
  \item \textbf{Probabilistic output.} Forecast intervals come built in.
\end{itemize}

It also has limits that bite hardest in finance:

\begin{itemize}
  \item \textbf{Weak signal.} Most of the variation in a monthly return is noise. A model trained to extend patterns may find them where none exist.
  \item \textbf{Domain mismatch.} Pretraining corpora are dominated by series with strong seasonality and trends: energy demand, traffic, sales. Financial returns have neither, and volatility has dynamics that few other series share.
  \item \textbf{Scale and context.} Normalizing each window discards the level of a series. For volatility or inflation, the level itself carries information about what comes next.
  \item \textbf{Contamination.}\index{contamination|textbf} Popular financial and macroeconomic series are widely published, so a model may have seen them, or close relatives, in pretraining.
\end{itemize}

Classical models are the opposite: narrow, but fitted to the target series and transparent about what they assume. The question is whether breadth beats fit.

\section{The question}
\label{sec:21-claim}

\begin{callout}
\textbf{The question.} Do pretrained time-series foundation models, used zero-shot, forecast financial and economic series better than classical models?
\end{callout}

To count as evidence, a test must compare against strong classical benchmarks for each series (the historical mean for returns, AR models for macroeconomic series, HAR for realized volatility\index{volatility!realized}), not only naive forecasts; use rolling origins with only data available at the time; report significance, not just average losses; and hold on data the model cannot have seen during pretraining.

\section{The original evidence}
\label{sec:21-evidence}

\textbf{The developers' benchmarks.} Chronos (Ansari et al.\ 2024), pretrained on public datasets and synthetic series, significantly outperformed other methods on datasets in its training corpus, and on new datasets was comparable to, and occasionally better than, models trained on them. Google reported that TimesFM (Das et al.\ 2024) beat ARIMA and ETS on unseen data and matched or beat deep-learning models trained on the target series. These benchmarks contain few financial series, and the Chronos result shows the \term{contamination} problem: performance was strongest on data the model had seen.

\textbf{Daily stock returns (Rahimikia, Ni and Wang 2025).} On daily excess returns\index{return!excess} in 94 countries, with an out-of-sample\index{in-sample versus out-of-sample} period from 2001 to 2023, zero-shot Chronos (large) had an out-of-sample $R^2$\index{out-of-sample $R^2$} of $-1.37\%$ and called the direction correctly about 51\% of the time, and TimesFM had $-2.80\%$. Fine-tuning helped little; training the same architectures on financial data helped substantially, but they still trailed tree ensembles such as CatBoost on $R^2$\index{regression!R-squared@$R^2$}.

\textbf{Multivariate financial series (Marconi 2025).} On Treasury yield changes, EUR/USD volatility and an equity spread, zero-shot models beat naive baselines on two series, and fine-tuning beat training from scratch with years less data, but specialized traditional models matched or beat the foundation model on two of three tasks.

\textbf{Realized volatility (Brini 2026).}\index{volatility!realized|textbf} Comparing nine zero-shot foundation models with eight econometric models on 50 assets at horizons of 1, 5 and 22 days, pooled losses favored the foundation models, but only because of a few unusual assets. On average per-asset QLIKE, only Tiny Time Mixers beat log-HAR at every horizon, by about 1.3\% to 1.8\%, and an equal-weight average of the two was among the best forecasts for almost every asset.

The pattern is consistent: zero-shot foundation models beat naive forecasts easily, sometimes match well-specified classical models, and rarely beat them by a meaningful margin. Where they do, combining them with a classical model helps more than replacing it.

\section{Lab~\thechapter: Zero-Shot Forecasts on Six Series}
\label{sec:21-lab}

\textbf{What you will explore:} how foundation models, used zero-shot, compare with classical forecasts. Parts~A, B and~D score the forecasts as in the zero-shot evaluations of Ansari et al.\ (2024) and Das et al.\ (2024) and the finance study of Rahimikia, Ni and Wang (2025), and Part~C checks whether a model's results depend on data it may have seen in training. The lab forecasts six free series one month ahead from every month since January 2000, with each model seeing only data up to that month: the U.S. market's excess return and its log realized variance (Kenneth R. French Data Library\index{Kenneth French data library}), and inflation, industrial production growth, and the changes in the unemployment rate and in the 10-year Treasury yield (FRED\index{FRED}). The classical models are the historical mean, the last value, an AR model chosen by BIC, which is \textbf{the benchmark}, a HAR-type model and GARCH(1,1); the foundation models are Chronos-Bolt and TimesFM, which are optional because they need PyTorch and large downloads.

\begin{itemize}
  \item \textbf{Part A} reports each model's MSE relative to AR and its $R^2$ against the historical mean (Lab Figures~\ref{ch:21}.1 and~\ref{ch:21}.5).
  \item \textbf{Part B} adds Diebold--Mariano tests, the share of months each model beats AR and the cumulative squared-error\index{cumulative squared-error plot} difference (Lab Figure~\ref{ch:21}.2).
  \item \textbf{Part C} splits the evaluation at each model's training cutoff, alongside a deliberately contaminated ``memorizing'' model (Lab Figure~\ref{ch:21}.3).
  \item \textbf{Part D} checks whether 80\% forecast intervals contain the outcome 80\% of the time (Lab Figure~\ref{ch:21}.4).
\end{itemize}

\label{sec:21-stress}
\textbf{Robustness checks.} Against the historical mean, AR and HAR rather than naive or untuned models, the gap narrows or disappears. Only data published after the training cutoff give a clean test\index{contamination|textbf}, and a year or two of monthly data gives tests little power. A few series with large errors can decide a pooled comparison, as Brini found, so report results per series and horizon, and name the exact model version and context length. When you have finished, compare what you saw with the published evidence summarized in Section~\ref{sec:21-verdict}.

\section{Where the evidence stands}
\label{sec:21-verdict}

The published studies point to an early picture: \textbf{zero-shot foundation models beat naive forecasts easily and sometimes match well-specified classical models, but they rarely beat them by a meaningful margin. The evidence is new and thin.} The table summarizes the evidence as of November 2026.

{\small
\begin{longtable}{@{}>{\raggedright\arraybackslash}p{0.19\linewidth}>{\raggedright\arraybackslash}p{0.13\linewidth}>{\raggedright\arraybackslash}p{\dimexpr 0.45\linewidth-6\tabcolsep\relax}>{\raggedright\arraybackslash}p{0.23\linewidth}@{}}
\toprule
Situation & Holds? & Why & Evidence \\
\midrule
\endfirsthead
\multicolumn{4}{@{}l}{\emph{(continued)}} \\
\toprule
Situation & Holds? & Why & Evidence \\
\midrule
\endhead
\midrule
\multicolumn{4}{r@{}}{\emph{(continued on the next page)}} \\
\endfoot
\bottomrule
\endlastfoot
The developers' benchmarks & \textbf{Partly} & Matched or beat dataset-specific models, most strongly on data in their training corpus, with few financial series. & Ansari et al.\ (2024); Das et al.\ (2024) \\*
\addlinespace[3pt]
Daily stock returns & \textbf{No} & Zero-shot out-of-sample $R^2$ of $-1.37\%$ and $-2.80\%$; financial pretraining helped. & Rahimikia, Ni and Wang (2025) \\
\addlinespace[3pt]
Realized volatility & \textbf{Weakly} & One of nine models beat log-HAR per asset, by 1 to 2\%; pooled wins came from a few outliers. & Brini (2026) \\
\addlinespace[3pt]
Other financial series & \textbf{Partly} & Beat naive baselines, but specialized models matched or beat them in two of three tasks. & Marconi (2025) \\
\addlinespace[3pt]
As part of a combination & \textbf{Partly} & An average with log-HAR was among the best forecasts for almost every asset. & Brini (2026) \\*
\addlinespace[3pt]
Macroeconomic series, after the cutoff & \textbf{Untested} & Clean post-cutoff evaluations are still scarce. & None yet \\
\end{longtable}
\addtocounter{table}{-1}% uncaptioned longtable: keep table numbers unchanged
}

The published evidence does not yet support a yes answer, but it is too new to rule one out. That financial pretraining helps suggests the problem is the domain, not the architecture.

\textbf{What would change the picture.} The answer would shift toward yes if a foundation model beat AR and HAR significantly on unseen data over several years, and toward no if successive generations kept falling short.

\section{What practitioners and investors can use}
\label{sec:21-investors}

Benchmark any foundation model against the historical mean, an autoregression and HAR, on the same rolling origins, and treat any evaluation before its training cutoff as possibly contaminated. Use foundation models as a second opinion, because averaging one with a classical model consistently beats either alone, and check that the intervals cover the outcome as often as they promise. Ask a vendor which classical benchmarks were used, what the training cutoff is, and how performance compares before and after it.

\begin{callout}
\textbf{What Failed: forecasting home prices at scale.} In 2018 Zillow began buying homes through Zillow Offers, relying on its price models to decide what to pay. In November 2021 it closed the business, wrote down its home inventory by \$304.4 million and cut about a quarter of its workforce, citing unpredictable home prices, capacity\index{capacity} constraints and other operational problems. A forecast that is accurate on average can still fail when acted on at scale, in a market that turns, against sellers who know more about the particular house.
\end{callout}

% Key terms of the chapter, in alphabetical order, for readers new to finance.
% A term defined in an earlier chapter points to the chapter, section and page.
\providecolor{FEMaroon}{RGB}{128,30,45}
\providecolor{FECream}{RGB}{253,248,238}
\providecolor{FEGold}{RGB}{196,150,60}
\makeatletter
\@ifundefined{keyterms}{%
  \newtcolorbox{keyterms}[1][Key terms in this chapter]{breakable, enhanced,
    colback=FECream, colframe=FEGold, boxrule=0.6pt, arc=2pt, leftrule=3pt,
    left=6pt, right=6pt, top=3pt, bottom=3pt, before skip=8pt, after skip=8pt,
    fonttitle=\bfseries, coltitle=white, colbacktitle=FEMaroon, title={#1}}%
}{}
\makeatother

\begin{keyterms}
\footnotesize
\begin{description}[leftmargin=0pt, itemsep=0pt, parsep=0pt, topsep=0pt, font=\normalfont\bfseries\color{FEMaroon}]
  \item[ARIMA.] ARIMA models extend an autoregression with differencing and moving-average terms.
  \item[Autoregression.] An autoregression forecasts a series by regressing it on its own past values.
  \item[Chronos.] Chronos is a time-series foundation model that scales each series, turns its values into tokens and forecasts the next token like a language model.
  \item[Contamination.] Contamination happens when the test data were part of a model's training data, which makes the model look better than it is.
  \item[Forecast interval calibration.] A forecast interval is well calibrated when the outcome falls inside it as often as promised, for example in 80\% of cases for an 80\% interval.
  \item[Foundation models.] Foundation models are large models trained on very broad collections of data and then used for many tasks with little or no further training.
  \item[Large language models.] Large language models are neural networks trained on vast amounts of text to predict the next word, which is how chatbots learn to write.
  \item[Pinball loss.] The pinball loss scores a forecast quantile, penalizing outcomes on the wrong side of it in proportion to the quantile's level.
  \item[Pretraining.] Pretraining is the first stage of building a foundation model, in which it learns general patterns from a very large collection of data before it is used on a specific task.
  \item[Rolling origins.] Rolling origins is the forecasting name for a walk-forward design, in which each forecast uses only data up to its own date.
  \item[TimesFM.] TimesFM is a time-series foundation model, a transformer that reads a series in patches and forecasts the next patch.
  \item[Training cutoff.] A model's training cutoff is the last date of the data it was trained on, so only later data give a clean test.
  \item[Transformer.] A transformer is the neural network design behind modern chatbots, built around the attention mechanism.
  \item[Zero-shot.] A model is used zero-shot when it forecasts a series without any training on that series.
\end{description}
\end{keyterms}
